% Zone „Das kennst du schon“ – aus bank/vierfeldertafel/zone.jsonl (zusammenbau v0.1)
\einheitenkopf[zone][Kennst du schon]{Das kennst du schon}
\setcounter{aufgabe}{0}

%% TODO Titel: Ich-kann-Satz fehlt in der Bank (Platzhalter: Kettenname) – Nr. 1
%% TODO Anweisung: ein Satz über der ersten Teilaufgabe fehlt in der Bank – Nr. 1
\begin{aufgabe}{Anteile, Prozentsätze und absolute Häufigkeiten ineinander umrechnen}
\begin{teile}
\teil $25\,\%$ von $80$? \leerfeld
\teil Anteil $0{,}35$ – wie viel Prozent, und wie viele von $240$? \leerfeld[\%] , \leerfeld
\end{teile}
\end{aufgabe}

%% TODO Titel: Ich-kann-Satz fehlt in der Bank (Platzhalter: Kettenname) – Nr. 2
%% TODO Anweisung: ein Satz über der ersten Teilaufgabe fehlt in der Bank – Nr. 2
\begin{aufgabe}{Zweistufige Zufallsexperimente und Pfadregeln (Baumdiagramm als Schwesterdarstellung, „Anteil vom Anteil“ als Produkt)}
\begin{teile}
\teil $\frac{1}{2}$ von $\frac{1}{3}$ – als Bruch? \leerfeld
\teil $30\,\%$ der Kunden sind Neukunden; $40\,\%$ der Neukunden kaufen Zubehör. Welcher Anteil aller Kunden sind Neukunden mit Zubehör? \leerfeld[\%]
\end{teile}
\end{aufgabe}

%% TODO Titel: Ich-kann-Satz fehlt in der Bank (Platzhalter: Kettenname) – Nr. 3
%% TODO Anweisung: ein Satz über der ersten Teilaufgabe fehlt in der Bank – Nr. 3
\begin{aufgabe}{Gegenereignis und Komplement (der Rest zu eins bzw. zur Gesamtzahl)}
\begin{teile}
\teil $P(A) = 0{,}3$ – $P(\overline{A})$? \leerfeld
\teil $P(A) = 0{,}125$ – $P(\overline{A})$? \leerfeld
\end{teile}
\end{aufgabe}

%% TODO Titel: Ich-kann-Satz fehlt in der Bank (Platzhalter: Kettenname) – Nr. 4
%% TODO Anweisung: ein Satz über der ersten Teilaufgabe fehlt in der Bank – Nr. 4
\begin{aufgabe}{Tabellen lesen und ausfüllen (Zeilen- und Spaltensummen)}
\begin{teile}
\teil Ergänze die fehlenden Zahlen der Tabelle.

\sachtabelle{lccc}{ & ja & nein & Summe}{Klasse 7a & 12 & \leerzelle & 28 \\ Klasse 7b & 15 & 10 & \leerzelle \\ Summe & \leerzelle & \leerzelle & \leerzelle}
\teil Eine Zeile einer Tabelle hat die Einträge $0{,}15$ und $0{,}37$. Zeilensumme? \leerfeld
\end{teile}
\end{aufgabe}

%% TODO Titel: Ich-kann-Satz fehlt in der Bank (Platzhalter: Kettenname) – Nr. 5
%% TODO Anweisung: ein Satz über der ersten Teilaufgabe fehlt in der Bank – Nr. 5
\begin{aufgabe}{Lineare Gleichungen und Ungleichungen mit einem Parameter lösen (negative Werte ausschließen)}
\begin{gleichungsraster}[2]
\gl{\text{$1 - 3p = 0{,}4$ – $p$?}} &
\end{gleichungsraster}
\begin{teile}
\teil Für welche $p > 0$ ist $1 - 2p - 3p$ nicht negativ? \leerfeld
\end{teile}
\end{aufgabe}

%% TODO Titel: Ich-kann-Satz fehlt in der Bank (Platzhalter: Kettenname) – Nr. 6
%% TODO Anweisung: ein Satz über der ersten Teilaufgabe fehlt in der Bank – Nr. 6
\begin{aufgabe}{Zweistufige Zufallsexperimente und Pfadregeln (Baumdiagramm als Schwesterdarstellung, „Anteil vom Anteil“ als Produkt) – Fehler finden}
\begin{teile}
\teil Ich finde den Fehler: $40\,\%$ der Mitglieder eines Vereins sind Jugendliche; $70\,\%$ der Jugendlichen spielen Fußball. Paul sagt: „$70\,\%$ aller Mitglieder sind fußballspielende Jugendliche.“
\end{teile}
\end{aufgabe}

%% TODO Titel: Ich-kann-Satz fehlt in der Bank (Platzhalter: Kettenname) – Nr. 7
%% TODO Anweisung: ein Satz über der ersten Teilaufgabe fehlt in der Bank – Nr. 7
\begin{aufgabe}{Zweistufige Zufallsexperimente und Pfadregeln (Baumdiagramm als Schwesterdarstellung, „Anteil vom Anteil“ als Produkt) – selbst rechnen}
\begin{teile}
\teil $80\,\%$ der Pakete gehen ins Inland; $15\,\%$ der Inlandspakete sind Expresspakete. Welcher Anteil aller Pakete sind Inland-Expresspakete? \leerfeld[\%]
\end{teile}
\end{aufgabe}
